% ChassisModule — 4輪統合足回り制御モジュール
\subsection{ChassisModule}

\begin{apibox}{ChassisModule --- 4輪統合足回り制御モジュール}
\textbf{定義ファイル:} \texttt{lib/ChassisModule/ChassisModule.h / .cpp}\\
\textbf{分類:} 機能統合モジュール（\texttt{updateOutput()}のみオーバーライド）\\
\textbf{対象デバイス:} 4輪駆動足回り（オムニ/メカナム/4WD）

内部で4つのDriveMotorModuleを保持し、前進・横移動・回転の3軸ベクトルを
ホイールパターンに応じた方向行列で各モーター出力に合成する統合モジュール。
\end{apibox}

\subsubsection{機能概要}

\begin{itemize}
    \item 4つのDriveMotorModuleの内部管理（生成・init・drive呼び出し）
    \item 3種類のホイールパターン対応（OMNI / MECANUM / FOURWD）
    \item 3軸ベクトル（前進・横移動・回転）から4輪出力への合成
    \item 出力リミッタ（\texttt{maxOutputPercent}で最大出力を制限）
    \item 各モーターの実出力値をデバッグ用に保持
\end{itemize}

% --- 定数 ---
\begin{funcblock}{定数}
\begin{tabularx}{\textwidth}{l l X}
\toprule
\textbf{名前} & \textbf{値} & \textbf{説明} \\
\midrule
\texttt{CHASSIS\_MOTOR\_COUNT} & \texttt{4} & モーター数（固定） \\
\bottomrule
\end{tabularx}
\end{funcblock}

% --- 列挙型 ---
\subsubsection{列挙型}

\begin{funcblock}{WheelPattern}
ホイール配置パターンを定義する。パターンによりベクトル合成の方向行列が異なる。

\begin{longtable}{l l p{6.0cm}}
\toprule
\textbf{値} & \textbf{名前} & \textbf{説明} \\
\midrule
\endhead
0 & \texttt{OMNI}    & オムニホイール。全方向移動 + 回転が可能 \\
1 & \texttt{MECANUM} & メカナムホイール。全方向移動 + 回転が可能（ローラー角度が異なる） \\
2 & \texttt{FOURWD}  & 4輪駆動。前後移動 + 回転のみ（横移動不可） \\
\bottomrule
\end{longtable}
\end{funcblock}

% --- Config ---
\subsubsection{Config構造体}

\begin{funcblock}{ChassisConfig}
4輪足回りの構成を定義する。

\begin{longtable}{l l l p{4.0cm}}
\toprule
\textbf{メンバ} & \textbf{型} & \textbf{制約・既定値} & \textbf{説明} \\
\midrule
\endhead
\texttt{motors[4]}        & \texttt{DriveMotorConfig[]} & ---         & 各モーターの設定。[0]:左前、[1]:右前、[2]:左後、[3]:右後 \\
\texttt{wheelPattern}     & \texttt{WheelPattern}       & ---         & ホイールパターン \\
\texttt{maxOutputPercent}  & \texttt{float}             & 例: 90.0    & 最大出力割合 [\%]。合成後の出力をこの値でスケーリング \\
\bottomrule
\end{longtable}
\end{funcblock}

% --- Data ---
\subsubsection{Data構造体}

\begin{funcblock}{ChassisData}
足回りの制御入力とデバッグ用出力を保持する。

\begin{longtable}{l l l p{4.0cm}}
\toprule
\textbf{メンバ} & \textbf{型} & \textbf{初期値} & \textbf{説明} \\
\midrule
\endfirsthead
\toprule
\textbf{メンバ} & \textbf{型} & \textbf{初期値} & \textbf{説明} \\
\midrule
\endhead
\bottomrule
\endlastfoot
\multicolumn{4}{l}{\textbf{--- 入力（ロジックフェーズで設定） ---}} \\
\texttt{forwardSpeed}  & \texttt{float} & \texttt{0.0f} & 前進速度 [$-100$\textasciitilde$100$]。正:前進、負:後退 \\
\texttt{lateralSpeed}  & \texttt{float} & \texttt{0.0f} & 横移動速度 [$-100$\textasciitilde$100$]。正:右、負:左 \\
\texttt{rotationSpeed} & \texttt{float} & \texttt{0.0f} & 回転速度 [$-100$\textasciitilde$100$]。正:右回転、負:左回転 \\
\midrule
\multicolumn{4}{l}{\textbf{--- 出力（updateOutput後に更新、デバッグ用） ---}} \\
\texttt{motorOutputs[4]} & \texttt{float[]} & \texttt{\{0\}} & 各モーターの実出力値 [\%]。[0]:左前、[1]:右前、[2]:左後、[3]:右後 \\
\end{longtable}
\end{funcblock}

% --- メソッド ---
\subsubsection{メソッド}

\begin{funcblock}{ChassisModule(const ChassisConfig\& config)}
\begin{tabularx}{\textwidth}{l X}
\toprule
\textbf{項目} & \textbf{内容} \\
\midrule
引数   & \texttt{config} --- \texttt{ChassisConfig}へのconst参照 \\
動作   & Configを\texttt{\_config}にコピー。4つのDriveMotorModuleを\texttt{config.motors[i]}から内部生成し、ポインタ配列\texttt{\_motors[]}にセット \\
\bottomrule
\end{tabularx}
\end{funcblock}

\begin{funcblock}{bool init()}
\begin{tabularx}{\textwidth}{l X}
\toprule
\textbf{項目} & \textbf{内容} \\
\midrule
戻り値   & \texttt{bool} --- 4モーター全ての\texttt{init()}が成功で\texttt{true}、1つでも失敗で\texttt{false} \\
動作     & 内部の4つのDriveMotorModuleの\texttt{init()}を順次呼び出す \\
\bottomrule
\end{tabularx}
\end{funcblock}

\begin{funcblock}{void updateOutput(SystemData\& data)}
\begin{tabularx}{\textwidth}{l X}
\toprule
\textbf{項目} & \textbf{内容} \\
\midrule
読み取り元     & \texttt{data.chassis} \\
タイミング制御 & 毎ループ実行 \\
動作           & ホイールパターンに応じた方向行列を取得し、\texttt{forwardSpeed}・\texttt{lateralSpeed}・\texttt{rotationSpeed}を合成して各モーター出力を計算。出力リミッタ適用後、各DriveMotorModuleの\texttt{drive()}を呼び出す。\texttt{motorOutputs[]}にデバッグ値を書き込む \\
\bottomrule
\end{tabularx}
\end{funcblock}

\begin{funcblock}{void deinit()}
\begin{tabularx}{\textwidth}{l X}
\toprule
\textbf{項目} & \textbf{内容} \\
\midrule
動作   & 全モーターを停止（\texttt{drive(0)}）させる \\
\bottomrule
\end{tabularx}
\end{funcblock}

% --- ベクトル合成 ---
\begin{notebox}[ベクトル合成と方向行列]
各モーターの出力は以下の式で計算される:

\vspace{0.3em}
$output_i = M_{i,0} \times forward + M_{i,1} \times lateral + M_{i,2} \times rotation$
\vspace{0.3em}

\noindent ここで $M$ はホイールパターンに応じた $4 \times 3$ の方向行列である。

\medskip
\textbf{OMNI（オムニホイール）の方向行列:}
\begin{center}
\begin{tabular}{l r r r}
\toprule
& \textbf{前進} & \textbf{横移動} & \textbf{回転} \\
\midrule
左前 [0] & $+1$ & $-1$ & $-1$ \\
右前 [1] & $+1$ & $+1$ & $+1$ \\
左後 [2] & $+1$ & $+1$ & $-1$ \\
右後 [3] & $+1$ & $-1$ & $+1$ \\
\bottomrule
\end{tabular}
\end{center}

\medskip
合成後の出力は絶対値の最大値が100を超える場合、全体を正規化して$\pm 100$以内に収める。
その後、\texttt{maxOutputPercent}でスケーリングする。
\end{notebox}

\begin{cautionbox}[DriveMotorModuleの管理]
ChassisModuleが内部で4つのDriveMotorModuleインスタンスを保持し、\texttt{init()}・\texttt{drive()}を呼び出す。
個別のDriveMotorModuleを\texttt{outputModules[]}に登録してはならない（二重制御になる）。
\end{cautionbox}

% --- 使用例 ---
\subsubsection{使用例}

\begin{lstlisting}[style=cppstyle, caption={ChassisModuleの使用例}]
// ProjectConfig.h
const ChassisConfig CHASSIS_CONFIG = {
    .motors = {
        {.in1Pin=4, .in2Pin=16, .pwmPin=17,
         .pwmChannel=2, .pwmFreqHz=1220,
         .minPowerThreshold=5.0f},  // 左前
        {.in1Pin=26, .in2Pin=25, .pwmPin=33,
         .pwmChannel=3, .pwmFreqHz=1220,
         .minPowerThreshold=5.0f},  // 右前
        {.in1Pin=18, .in2Pin=19, .pwmPin=23,
         .pwmChannel=4, .pwmFreqHz=1220,
         .minPowerThreshold=5.0f},  // 左後
        {.in1Pin=13, .in2Pin=14, .pwmPin=27,
         .pwmChannel=5, .pwmFreqHz=1220,
         .minPowerThreshold=5.0f},  // 右後
    },
    .wheelPattern     = WheelPattern::OMNI,
    .maxOutputPercent = 90.0f,
};

// ロジックフェーズでの使用
void applyPattern(SystemData& data) {
    // BLE受信データから操作指令を反映
    if (data.ble.newDataIn) {
        data.chassis.forwardSpeed  = (int8_t)data.ble.rxData[0];
        data.chassis.lateralSpeed  = (int8_t)data.ble.rxData[1];
        data.chassis.rotationSpeed = (int8_t)data.ble.rxData[2];
        data.ble.newDataIn = false;
    }
}
\end{lstlisting}
